% 19-limitations.tex: current limitations and what external validation is required.
\section{Current limitations and the external validation still required}
\label{sec:limitations}

The repository's readiness ledger opens with one line, and this paper repeats it as its own bound:
at \code{v9.345} Polaris is not production-ready for real identity data. Every engineering gap the
ledger enumerated is closed and pinned by a check; what remains is not buildable in the repository.
The list below separates what has not been done from what cannot be done by the author.

\subsection{Decisions only a deploying organisation can make}

Nine decisions must be recorded for a named deployment before real data could be held: a legal
basis, impact assessment and regulator approval; the signing-key custody driver and who holds the
key; the high-availability topology and its hosts; encryption at rest and its key custodian; the
off-site backup target and schedule; the alerting backend, the on-call rotation and who receives the
duress page; the right-to-erasure policy; the retention schedule per class, with the counsel who says
the number satisfies the statute; and an independent penetration test with a threat-model sign-off.
Polaris records each decision and its justification; it does not know the law. \sExt

\subsection{What is not built}

\begin{itemize}
  \item The physical token is modelled, not manufactured; the enrolment kit, the card profile and
    the verifier device are roadmap rows.
  \item No hardware security module has been used; the PKCS\#11 driver is proven against a software
    token, and the lattice library is not FIPS-validated.
  \item The profile is single-region; multi-region recovery, status distribution at content-delivery
    scale, and the epoch pipeline at national scale are not built.
  \item A hash-based signer is registered and not wired. Two internal transport hops, the
    certificate signatures and today's operator authenticators remain classical.
  \item Blinded and one-time presentations are not built, so a verifier looking at a full
    credential still sees a stable token value, an issuer signature and a holder key. The
    per-relying-party handle bounds what a verifier \emph{stores}, not what it is \emph{shown};
    Section~\ref{sec:privacy} states the distinction and the verifier reports which of the two a
    given presentation gave. This is the residue of the holder-side work, not the whole of it: the
    holder key, on-device proving, the scoped nullifier and the pairwise handle were built at
    \code{v9.349} through \code{v9.353} and this list said otherwise until \code{v9.355}.
  \item Neither the scoped nullifier nor the pairwise handle hides a holder from the \emph{issuer},
    which derives every epoch leaf in order to build the tree and could therefore compute any
    member's nullifier in any scope. Issuer-blind derivation is a different construction and is not
    claimed.
  \item The interim credential during a recovery cool-down is not built.
\end{itemize}

\subsection{What has not been validated by anyone else}

\begin{itemize}
  \item No external penetration test, cryptographic audit, red team, bug bounty or certification
    has occurred. The scope document for the engagement exists; no engagement has been commissioned.
  \item No pilot with consenting users has run; the data has always been notional.
  \item No external team has implemented the protocol from the specification and conformance suite
    alone; the suite is the contract and the integration has not happened.
  \item No institutionally independent witness watches any log; the witnesses in the record are
    processes in a drill. Two federated instances are two instances of one code base run by one
    author.
  \item The circuit over the proof library has had no external review, and the concrete bit
    security of the shipped proof configuration has not been independently derived.
  \item The duress path is a tested property of the modelled flow, not an audited side-channel
    guarantee; long-run frequency analysis is an accepted residual.
  \item The repository's own thesis, that a stranger can read it cold and correctly infer what it
    needs, was never tested by a stranger before its deadline and is recorded as inconclusive, with
    the strong claim retired.
\end{itemize}

\subsection{What the numbers are}

Every performance number in this paper is single-node on one laptop, or a simulation at a stated
scale factor, or a projection that the repository labels as one. No production cluster has been
measured. The national planning targets are targets.

\subsection{Documentation that lags the code}

Researching this version found a few places where the repository's own prose is behind its code
or its history; they are recorded in the author's notes that accompany this version rather than
fixed silently. Where this paper and any repository document disagree, the paper followed the code
at \code{v9.345} and says so in the note.
